\chapter{ER-to-Relational Mapping}

\section{Mapping Rules}
The standard rules were applied, with one project-wide convention: every entity gets a surrogate integer
primary key (\tbl{AUTO\_INCREMENT}), and natural identifiers become UNIQUE keys. Names, phone numbers and batch
numbers are not reliable identifiers, and narrow integer keys keep foreign keys and InnoDB's clustered index small.
\begin{enumerate}
\item \textbf{Strong entity} $\to$ a table with its simple attributes; its key becomes the primary key.
\item \textbf{Weak entity} $\to$ a table holding the owner's key as a foreign key (NOT NULL) plus the partial key; (owner key, partial key) is a UNIQUE key. A surrogate primary key is added.
\item \textbf{1:1} $\to$ the foreign key goes on the side that may be missing, with a UNIQUE constraint.
\item \textbf{1:N} $\to$ the foreign key goes on the N side; NOT NULL when the N side's participation is total.
\item \textbf{M:N and ternary} $\to$ a new table with a foreign key to each participant and the relationship's attributes.
\item \textbf{Composite attribute} $\to$ its components become columns. \textbf{Multivalued attribute} $\to$ a separate table (owner key, value), or (as here, Section~4.2) one column when the value is never searched element by element. \textbf{Derived attribute} $\to$ not stored, or a generated column, or a stored copy kept up to date.
\end{enumerate}

\section{Step-by-Step Conversion}
\begin{tabularx}{\textwidth}{@{}R{3.1cm}Y@{}}
\toprule
\textbf{ER construct} & \textbf{Relation in hospital\_db (real example)} \\\midrule
Strong entity & PATIENT $\to$ \tbl{patients(\underline{patient\_id}, name, date\_of\_birth, gender, \dots, abha\_number)}; ABHA becomes UNIQUE \tbl{uq\_patients\_abha}. Likewise for the 25 other strong entities. \\
Weak entity & BED (owner WARD) $\to$ \tbl{beds(\underline{bed\_id}, ward\_id\textsuperscript{FK}, bed\_number, \dots)} with UNIQUE (\tbl{ward\_id, bed\_number}); MEDICINE\_BATCH $\to$ \tbl{medicine\_batches} with UNIQUE (\tbl{medicine\_id, batch\_number}); DOCTOR\_SCHEDULE and DAILY\_ANALYTICS likewise, with ON DELETE CASCADE from \tbl{doctors}. \\
1:1 & \textsc{becomes}: \tbl{encounters.appointment\_id} (NULL for walk-ins) with UNIQUE \tbl{uq\_encounters\_appointment}; \textsc{is\_doctor}: \tbl{doctors.user\_id} NOT NULL UNIQUE (\tbl{uq\_doctors\_user\_id}); \textsc{yields}: \tbl{lab\_results.order\_id} NOT NULL UNIQUE. \\
1:N & \textsc{books}: \tbl{appointments.patient\_id} NOT NULL (total); \textsc{carries}: \tbl{emergency\_cases.ambulance\_id} NULL (partial). 64 relationships are mapped this way. \\
M:N with attribute & \textsc{filled\_from} $\to$ \tbl{prescription\_item\_batches(\underline{id}, item\_id\textsuperscript{FK}, batch\_id\textsuperscript{FK}, quantity)}, UNIQUE (\tbl{item\_id, batch\_id}). \\
Ternary with attributes & \textsc{occupies} $\to$ \tbl{bed\_assignment\_log(\underline{log\_id}, bed\_id\textsuperscript{FK}, patient\_id\textsuperscript{FK}, admission\_id\textsuperscript{FK}, assigned\_at, released\_at, status)}. \\
Composite & \emph{blood\_pressure} $\to$ \tbl{bp\_systolic, bp\_diastolic} (CHECK: both or neither, systolic > diastolic); \emph{slot} $\to$ \tbl{appointment\_date, appointment\_time}. \\
Multivalued & \emph{allergies} is stored as one text column \tbl{patients.allergies}: the application only displays it on the chart and never queries one allergy, so a separate \tbl{patient\_allergies} table was not created (a limitation, Section~11.2). \\
Derived & \emph{age}: not stored, computed as \tbl{TIMESTAMPDIFF(YEAR, date\_of\_birth, ?)} with today's date in India as the parameter; \emph{total\_price}: generated column \tbl{quantity * unit\_price}; \emph{stock\_quantity}, \emph{total\_amount}, \emph{current\_workload}: stored copies kept equal by transactions and triggers (Section~6.6). \\
\bottomrule
\end{tabularx}

\section{Final Relational Schema}
All 33 relations, with columns in table order. \underline{Underlined}: primary key. \textit{Italic}\textsuperscript{FK}:
foreign key (with its target when the name does not say it). \textsuperscript{G}: generated column. \textsc{uk}: other
candidate keys (UNIQUE). Not shown: the 15 composite consistency keys of Section~6.5, which re-check pairs of
these foreign keys.

\begin{footnotesize}
\relation{patients}{\underline{patient\_id}, name, date\_of\_birth, dob\_estimated, gender, blood\_group, phone, address, emergency\_contact, created\_at, updated\_at, allergies, abha\_number}{(abha\_number)}
\relation{appointments}{\underline{appointment\_id}, \fk{patient\_id}{}, \fk{doctor\_id}{}, appointment\_date, appointment\_time, status, reason, notes, created\_at, updated\_at}{(doctor\_id, appointment\_date, appointment\_time)}
\relation{encounters}{\underline{encounter\_id}, \fk{patient\_id}{}, \fk{doctor\_id}{}, \fk{appointment\_id}{}, \fk{admission\_id}{}, encounter\_type, status, arrived\_at, triaged\_at, start\_timestamp, end\_timestamp, \fk{created\_by}{users}, backfilled\_from, created\_at, updated\_at}{(appointment\_id); (backfilled\_from)}
\relation{consultations}{\underline{consultation\_id}, \fk{appointment\_id}{}, \fk{patient\_id}{}, \fk{doctor\_id}{}, \fk{admission\_id}{}, symptoms, diagnosis, assessment, plan, notes, consultation\_time, created\_at, \fk{encounter\_id}{}}{}
\relation{vital\_signs}{\underline{vital\_id}, \fk{patient\_id}{}, \fk{encounter\_id}{}, \fk{admission\_id}{}, recorded\_at\_utc, \fk{recorded\_by}{users}, temperature\_c, pulse\_bpm, resp\_rate, bp\_systolic, bp\_diastolic, spo2\_pct, pain\_score, notes, created\_at}{}
\relation{surgery\_requests}{\underline{request\_id}, \fk{patient\_id}{}, \fk{doctor\_id}{}, procedure\_name, diagnosis, priority, requested\_date, status, notes, created\_at, updated\_at, \fk{encounter\_id}{}}{}
\relation{departments}{\underline{department\_id}, name, description, floor, phone, is\_active}{(name)}
\relation{wards}{\underline{ward\_id}, name, \fk{department\_id}{}, floor, capacity, ward\_type}{}
\relation{beds}{\underline{bed\_id}, \fk{ward\_id}{}, bed\_number, bed\_type, status, has\_ventilator}{(ward\_id, bed\_number)}
\relation{admissions}{\underline{admission\_id}, \fk{patient\_id}{}, \fk{doctor\_id}{}, \fk{bed\_id}{}, \fk{department\_id}{}, \fk{emergency\_id}{}, admission\_date, discharge\_date, status, diagnosis, notes, created\_at, updated\_at}{}
\relation{bed\_assignment\_log}{\underline{log\_id}, \fk{bed\_id}{}, \fk{patient\_id}{}, \fk{admission\_id}{}, assigned\_at, released\_at, status}{}
\relation{emergency\_cases}{\underline{emergency\_id}, \fk{patient\_id}{}, \fk{ambulance\_id}{}, severity, symptoms, arrival\_time, required\_specialization, required\_bed\_type, ventilator\_required, status, \fk{assigned\_doctor\_id}{doctors}, \fk{assigned\_bed\_id}{beds}, notes, created\_at, updated\_at}{}
\relation{ambulances}{\underline{ambulance\_id}, vehicle\_number, driver\_name, driver\_phone, current\_location, status}{(vehicle\_number)}
\relation{medicines}{\underline{medicine\_id}, name, category, manufacturer, stock\_quantity, reorder\_level, unit\_price, expiry\_date, requires\_prescription, created\_at, updated\_at}{}
\relation{medicine\_batches}{\underline{batch\_id}, \fk{medicine\_id}{}, batch\_number, quantity, expiry\_date, received\_date, \fk{received\_by}{users}, created\_at, updated\_at}{(medicine\_id, batch\_number)}
\relation{prescriptions}{\underline{prescription\_id}, \fk{consultation\_id}{}, \fk{patient\_id}{}, \fk{doctor\_id}{}, prescription\_date, status, notes, created\_at, \fk{encounter\_id}{}, \fk{admission\_id}{}}{}
\relation{prescription\_items}{\underline{item\_id}, \fk{prescription\_id}{}, \fk{medicine\_id}{}, dosage, frequency, duration, quantity, dispensed, dispensed\_at, frequency\_code, duration\_days, route, units\_per\_dose}{}
\relation{prescription\_item\_batches}{\underline{id}, \fk{item\_id}{}, \fk{batch\_id}{}, quantity, created\_at}{(item\_id, batch\_id)}
\relation{pharmacy\_stock\_movements}{\underline{movement\_id}, \fk{medicine\_id}{}, movement\_type, quantity, reason, notes, \fk{performed\_by}{users}, created\_at, reference\_id, \fk{batch\_id}{}}{}
\relation{medication\_administrations}{\underline{administration\_id}, \fk{prescription\_item\_id}{prescription\_items}, \fk{admission\_id}{}, scheduled\_at\_utc, status, administered\_at\_utc, dose\_given, route\_given, reason, \fk{recorded\_by}{users}, recorded\_at\_utc, created\_at, updated\_at}{(prescription\_item\_id, scheduled\_at\_utc)}
\relation{lab\_tests}{\underline{test\_id}, name, category, \fk{department\_id}{}, normal\_range, unit, cost, reference\_low, reference\_high, critical\_low, critical\_high}{}
\relation{lab\_orders}{\underline{order\_id}, \fk{consultation\_id}{}, \fk{patient\_id}{}, \fk{doctor\_id}{}, \fk{test\_id}{}, order\_date, status, priority, notes, created\_at, updated\_at, \fk{encounter\_id}{}, \fk{admission\_id}{}}{}
\relation{lab\_results}{\underline{result\_id}, \fk{order\_id}{}, result\_value, unit, reference\_range, interpretation, result\_date, technician\_notes, \fk{performed\_by}{users}, reviewed\_by\_doctor, reviewed\_at, created\_at, numeric\_value, reference\_low, reference\_high, interpretation\_source}{(order\_id)}
\relation{lab\_result\_attachments}{\underline{attachment\_id}, \fk{result\_id}{}, original\_filename, stored\_filename, mime\_type, file\_size, \fk{uploaded\_by}{users}, uploaded\_at}{(stored\_filename)}
\relation{users}{\underline{user\_id}, username, password\_hash, role, full\_name, email, phone, is\_active, created\_at, updated\_at}{(email); (username)}
\relation{doctors}{\underline{doctor\_id}, \fk{user\_id}{}, \fk{department\_id}{}, name, specialization, phone, status, shift, current\_workload, max\_workload, created\_at, updated\_at}{(user\_id)}
\relation{doctor\_schedules}{\underline{schedule\_id}, \fk{doctor\_id}{}, day\_of\_week, start\_time, end\_time, is\_available}{(doctor\_id, day\_of\_week, start\_time)}
\relation{doctor\_daily\_analytics}{\underline{analytics\_id}, \fk{doctor\_id}{}, analytics\_date, patients\_seen, emergency\_cases, appointment\_patients, patients\_treated, cases\_solved, cases\_postponed, cases\_cancelled, cases\_pending, created\_at, updated\_at}{(doctor\_id, analytics\_date)}
\relation{nurse\_ward\_assignments}{\underline{assignment\_id}, \fk{nurse\_user\_id}{users}, \fk{ward\_id}{}, start\_date, end\_date, \fk{assigned\_by}{users}, created\_at, updated\_at}{}
\relation{bills}{\underline{bill\_id}, \fk{patient\_id}{}, \fk{admission\_id}{}, bill\_date, total\_amount, paid\_amount, status, payment\_method, notes, created\_at, updated\_at}{}
\relation{bill\_items}{\underline{item\_id}, \fk{bill\_id}{}, description, category, quantity, unit\_price, total\_price\gen, created\_at}{}
\relation{audit\_logs}{\underline{audit\_id}, \fk{user\_id}{}, role, action, entity\_type, entity\_id, patient\_id, outcome, status\_code, method, path, ip\_address, details, created\_at}{}
\relation{schema\_migrations}{\underline{version}, applied\_at, note}{}

\end{footnotesize}
